% GENERATED FILE. Edit canonical lessons, then run npm run generate:books.
% canonical-source-hash: fnv1a64:359d7c64d34171ee
% canonical-lessons: IT-C378-scopo, IT-C378-intenzione, IT-C378-programmare, IT-C378-proposito, IT-C378-meta, IT-R378-ripasso-371-378

\chapter{Che cosa voglio fare}
\label{ch:it-che-cosa-voglio-fare}
\hlchaptermodality{378}

\begin{chapteropening}
\textbf{By the end of this chapter:} \emph{I can say what I aim to do: my purpose, my intention, a plan, a resolution and a goal.}

Complete the last lesson of chapter 378: I can say what I aim to do: my purpose, my intention, a plan, a resolution and a goal.
\end{chapteropening}

\section[lo scopo]{lo scopo — the purpose, the aim}

\label{lesson:IT-C378-scopo}



\begin{quote}
Before the new one: say the Italian for to commit yourself, then the Italian for flexible.
\end{quote}

\subsection*{You'll want to know: lo scopo}
\textbf{lo scopo} — \textquotedblleft{}the purpose, the aim\textquotedblright{}.

\begin{quote}
Lo scopo del viaggio è imparare la lingua.
\end{quote}

\begin{quote}
\emph{The aim of the trip is to learn the language.}
\end{quote}

\emph{Allo scopo di} — in order to. \emph{A che scopo?} — what for?

\begin{cousinweb}
\emph{scopo} is a learned borrowing from Greek \emph{skopós} \textquotedblleft{}a watcher; a mark to aim at\textquotedblright{}, through Latin \emph{scopos} \textquotedblleft{}a mark to aim at\textquotedblright{} (Lewis \& Short), which EtymDB gives as the source of \emph{scopo} (EtymDB); it is related to \emph{sképtesthai} \textquotedblleft{}to look carefully\textquotedblright{}.
\end{cousinweb}

\begin{grammarlens}[title={What you've built}]
One more word for telling, explaining and describing at B1.
\end{grammarlens}

\subsection*{Your turn}
\begin{itemize}
\raggedright
  \item \emph{Say it:} \emph{lo scopo}
  \item \emph{Say it:} \emph{lo scopo}, once more
  \item \emph{Say it:} a sentence of your own with \emph{lo scopo}
  \item \emph{Recall:} say the Italian for to commit yourself, then the Italian for flexible, then say \emph{lo scopo} again
\end{itemize}

\begin{tcolorbox}[breakable,colback=teal!4,colframe=teal!35!black,title={Before you move on}]
What does lo scopo mean? (\textquotedblleft{}The purpose, the aim\textquotedblright{}.) Say it once more.
\end{tcolorbox}
\section[l'intenzione]{l'intenzione — the intention}

\label{lesson:IT-C378-intenzione}



\begin{quote}
Before the new one: say the Italian for flexible, then the Italian for the purpose.
\end{quote}

\subsection*{You'll want to know: l'intenzione}
\textbf{l'intenzione} — \textquotedblleft{}the intention\textquotedblright{}.

\begin{quote}
Non ho intenzione di cambiare lavoro.
\end{quote}

\begin{quote}
\emph{I have no intention of changing jobs.}
\end{quote}

\emph{Avere intenzione di} + infinitive — to intend to.

\begin{cousinweb}
\emph{intenzione} is a learned borrowing from Latin \emph{intentiō} \textquotedblleft{}a stretching, an effort, a purpose\textquotedblright{} (Lewis \& Short), from \emph{intendere} \textquotedblleft{}to stretch towards\textquotedblright{}, built on \emph{tendere} \textquotedblleft{}to stretch\textquotedblright{}.
\end{cousinweb}

\begin{grammarlens}[title={What you've built}]
One more word for telling, explaining and describing at B1.
\end{grammarlens}

\subsection*{Your turn}
\begin{itemize}
\raggedright
  \item \emph{Say it:} \emph{l'intenzione}
  \item \emph{Say it:} \emph{l'intenzione}, once more
  \item \emph{Say it:} a sentence of your own with \emph{l'intenzione}
  \item \emph{Recall:} say the Italian for flexible, then the Italian for the purpose, then say \emph{l'intenzione} again
\end{itemize}

\begin{tcolorbox}[breakable,colback=teal!4,colframe=teal!35!black,title={Before you move on}]
What does l'intenzione mean? (\textquotedblleft{}The intention\textquotedblright{}.) Say it once more.
\end{tcolorbox}
\section[programmare]{programmare — to plan; to program}

\label{lesson:IT-C378-programmare}



\begin{quote}
Before the new one: say the Italian for the purpose, then the Italian for the intention.
\end{quote}

\subsection*{You'll want to know: programmare}
\textbf{programmare} — \textquotedblleft{}to plan; to program\textquotedblright{}.

\begin{quote}
Abbiamo programmato le vacanze con sei mesi di anticipo.
\end{quote}

\begin{quote}
\emph{We planned the holidays six months in advance.}
\end{quote}

Also to program a computer or a machine.

\begin{cousinweb}
\emph{programmare} is built on \emph{programma} \textquotedblleft{}programme\textquotedblright{}, a learned borrowing from Late Latin \emph{programma} \textquotedblleft{}a proclamation\textquotedblright{} (EtymDB), from Greek \emph{prógramma} \textquotedblleft{}a public notice\textquotedblright{}, from \emph{prográphein} \textquotedblleft{}to write before\textquotedblright{}, on \emph{gráphein} \textquotedblleft{}to write\textquotedblright{}.
\end{cousinweb}

\begin{grammarlens}[title={What you've built}]
One more word for telling, explaining and describing at B1.
\end{grammarlens}

\subsection*{Your turn}
\begin{itemize}
\raggedright
  \item \emph{Say it:} \emph{programmare}
  \item \emph{Say it:} \emph{programmare}, once more
  \item \emph{Say it:} a sentence of your own with \emph{programmare}
  \item \emph{Recall:} say the Italian for the purpose, then the Italian for the intention, then say \emph{programmare} again
\end{itemize}

\begin{tcolorbox}[breakable,colback=teal!4,colframe=teal!35!black,title={Before you move on}]
What does programmare mean? (\textquotedblleft{}To plan; to program\textquotedblright{}.) Say it once more.
\end{tcolorbox}
\section[il proposito]{il proposito — the resolution, the intention}

\label{lesson:IT-C378-proposito}



\begin{quote}
Before the new one: say the Italian for the intention, then the Italian for to plan.
\end{quote}

\subsection*{You'll want to know: il proposito}
\textbf{il proposito} — \textquotedblleft{}the resolution, the intention\textquotedblright{}.

\begin{quote}
Il mio proposito per l'anno nuovo è fare più sport.
\end{quote}

\begin{quote}
\emph{My New Year's resolution is to do more sport.}
\end{quote}

\emph{A proposito!} — by the way! \emph{Di proposito} — on purpose.

\begin{cousinweb}
\emph{proposito} is a learned borrowing from Latin \emph{prōpositum} \textquotedblleft{}a plan, a purpose\textquotedblright{} (Lewis \& Short), from \emph{prōpōnere} \textquotedblleft{}to set forth\textquotedblright{}, built on \emph{pōnere} \textquotedblleft{}to put\textquotedblright{}.
\end{cousinweb}

\begin{grammarlens}[title={What you've built}]
One more word for telling, explaining and describing at B1.
\end{grammarlens}

\subsection*{Your turn}
\begin{itemize}
\raggedright
  \item \emph{Say it:} \emph{il proposito}
  \item \emph{Say it:} \emph{il proposito}, once more
  \item \emph{Say it:} a sentence of your own with \emph{il proposito}
  \item \emph{Recall:} say the Italian for the intention, then the Italian for to plan, then say \emph{il proposito} again
\end{itemize}

\begin{tcolorbox}[breakable,colback=teal!4,colframe=teal!35!black,title={Before you move on}]
What does il proposito mean? (\textquotedblleft{}The resolution, the intention\textquotedblright{}.) Say it once more.
\end{tcolorbox}
\section[la meta]{la meta — the goal; the destination}

\label{lesson:IT-C378-meta}



\begin{quote}
Before the new one: say the Italian for to plan, then the Italian for the resolution.
\end{quote}

\subsection*{You'll want to know: la meta}
\textbf{la meta} — \textquotedblleft{}the goal; the destination\textquotedblright{}.

\begin{quote}
La nostra meta è la cima della montagna.
\end{quote}

\begin{quote}
\emph{Our goal is the top of the mountain.}
\end{quote}

Stress on the first syllable: \emph{mèta}. Not to be mixed up with \emph{metà} \textquotedblleft{}half\textquotedblright{}.

\begin{cousinweb}
\emph{meta} is a learned borrowing from Latin \emph{mēta} \textquotedblleft{}a cone-shaped turning post in a racecourse; a limit, a goal\textquotedblright{} (Lewis \& Short).
\end{cousinweb}

\begin{grammarlens}[title={What you've built}]
One more word for telling, explaining and describing at B1.
\end{grammarlens}

\subsection*{Your turn}
\begin{itemize}
\raggedright
  \item \emph{Say it:} \emph{la meta}
  \item \emph{Say it:} \emph{la meta}, once more
  \item \emph{Say it:} a sentence of your own with \emph{la meta}
  \item \emph{Recall:} say the Italian for to plan, then the Italian for the resolution, then say \emph{la meta} again
\end{itemize}

\begin{tcolorbox}[breakable,colback=teal!4,colframe=teal!35!black,title={Before you move on}]
What does la meta mean? (\textquotedblleft{}The goal; the destination\textquotedblright{}.) Say it once more.
\end{tcolorbox}
\section[(dialogue)]{Ripasso dei capitoli 371–378}

\label{lesson:IT-R378-ripasso-371-378}



\begin{quote}
Say the Italian for relevant and crucial.
\end{quote}

\subsection*{Your turn}
Say these aloud:

\begin{itemize}
\raggedright
  \item relevant, crucial, essential, fundamental, the premise
  \item to suppose, the hypothesis, to presume, plausible, the precondition
  \item provided that, if, otherwise, or else, possible
  \item the unexpected problem, predictable, unpredictable, the alternative, to adapt
  \item to cancel, to postpone, provisional, final, to fix
  \item the risk, risky, likely, unlikely, the probability
  \item available, busy, the commitment, to commit yourself, flexible
  \item the purpose, the intention, to plan, the resolution, the goal
\end{itemize}

\begin{tcolorbox}[breakable,colback=teal!4,colframe=teal!35!black,title={Before you move on}]
Relevant? (\textbf{rilevante}.) Crucial? (\textbf{cruciale}.) Essential? (\textbf{essenziale}.) Fundamental? (\textbf{fondamentale}.) The premise? (\textbf{la premessa}.) To suppose? (\textbf{supporre}.) The hypothesis? (\textbf{l'ipotesi}.) To presume? (\textbf{presumere}.) Plausible? (\textbf{plausibile}.) The precondition? (\textbf{il presupposto}.) Provided that? (\textbf{purché}.) If? (\textbf{qualora}.) Otherwise? (\textbf{altrimenti}.) Or else? (\textbf{sennò}.) Possible? (\textbf{eventuale}.) The unexpected problem? (\textbf{l'imprevisto}.) Predictable? (\textbf{prevedibile}.) Unpredictable? (\textbf{imprevedibile}.) The alternative? (\textbf{l'alternativa}.) To adapt? (\textbf{adattarsi}.) To cancel? (\textbf{annullare}.) To postpone? (\textbf{rinviare}.) Provisional? (\textbf{provvisorio}.) Final? (\textbf{definitivo}.) To fix? (\textbf{stabilire}.) The risk? (\textbf{il rischio}.) Risky? (\textbf{rischioso}.) Likely? (\textbf{probabile}.) Unlikely? (\textbf{improbabile}.) The probability? (\textbf{la probabilità}.) Available? (\textbf{disponibile}.) Busy? (\textbf{impegnato}.) The commitment? (\textbf{l'impegno}.) To commit yourself? (\textbf{impegnarsi}.) Flexible? (\textbf{flessibile}.) The purpose? (\textbf{lo scopo}.) The intention? (\textbf{l'intenzione}.) To plan? (\textbf{programmare}.) The resolution? (\textbf{il proposito}.) The goal? (\textbf{la meta}.)
\end{tcolorbox}
